Source : OCDE, Taxing Wages, indicateurs comparatifs, données 2025.\par Lecture : pour 100 \euro{} que coûte à son employeur un salarié célibataire sans enfant payé au salaire moyen, 47 \euro{} vont en France aux cotisations et à l'impôt sur le revenu, dont 27 \euro{} de cotisations employeur ; le salarié en garde 53. C'est un taux moyen, au salaire moyen : il est plus élevé pour un salaire plus élevé. L'OCDE publie l'impôt et les cotisations en pourcentage du salaire brut ; ils sont ici rapportés au coût du travail, pour que les trois composantes s'additionnent.
